\documentclass[12pt]{article}
\usepackage[utf8]{inputenc}
\usepackage[T1]{fontenc}
\usepackage{hyperref}
\title{Tasmanian, Oyster Bay Language Resources for K-12 Curriculum Development}
\author{Compiled by Claude Sonnet 5.5 (claude-sonnet-5-5) for Andr\'as Kornai}
\date{}
\begin{document}
\maketitle

\section{Introduction}
Oyster Bay Tasmanian (Paredarerme; ISO 639-3 xpd, Glottocode oyst1235) was spoken on Tasmania's central east coast. It is extinct and is treated as an isolate. Its documentation is better than that of most Tasmanian varieties. Robinson recorded a list of 357 words, and Joseph Milligan published a coastal list of 1,040 words in 1857 and 1859, the longest vocabulary of any Tasmanian variety. Little Swanport is sometimes considered a dialect. Oyster Bay data are also a basis of palawa kani, the Tasmanian Aboriginal Centre's revival language. Two Oyster Bay lists reportedly contain substantial northeastern admixture. The searches did not return the word lists themselves.

\section{Primary Educational Resources}
No primers, textbooks or school programmes for Oyster Bay itself were found. The palawa kani revival is described at: \url{https://en.wikipedia.org/wiki/Palawa_kani}

\section{Linguistic and Reference Materials}
\begin{itemize}
\item Paredarerme language (Wikipedia): \url{https://en.wikipedia.org/wiki/Paredarerme_language}
\item Little Swanport language (Wikipedia): \url{https://en.wikipedia.org/wiki/Little_Swanport_language}
\item Tasmanian languages (Wikipedia): \url{https://en.wikipedia.org/wiki/Tasmanian_languages}
\item Eastern Tasmanian languages (Wikipedia): \url{https://en.wikipedia.org/wiki/Eastern_Tasmanian_languages}
\item AIATSIS AUSTLANG entry T2, where the original sources can be traced: \url{https://collection.aiatsis.gov.au/austlang/language/T2}
\item ATSIDA archive record, returned by the search and not examined: \url{https://www.atsida.edu.au/node/363}
\end{itemize}

\section{Religious and Community Resources}
The Tasmanian Aboriginal Centre runs palawa kani. The searches indicate it restricts access to its learning materials to the Aboriginal community and claims copyright in them. Any use requires its permission.

\section{Digital and Audio Resources}
No public audio, video or software was found.

\section{Curriculum Development Context}
The language has no speakers, but its historical vocabulary is relatively rich. A school programme should be developed with the Tasmanian Aboriginal community. It could present the history of the languages, the nature of the Robinson and Milligan records, and the palawa kani revival, and it should not teach Oyster Bay vocabulary independently of the community. The recorders used English-based spellings, so the original sounds are uncertain.
\end{document}